\title{QLoRA: Efficient Finetuning of Quantized LLMs}

\begin{abstract}
We introduce \textsc{QLoRA}, a memory-efficient finetuning method that enables the adaptation of large-scale models with up to 65 billion parameters on a single 48GB GPU, while maintaining the effectiveness typically seen in full 16-bit finetuning. \textsc{QLoRA} facilitates gradient propagation through a fixed, 4-bit quantized pretrained language model into Low Rank Adapters~(LoRA). Our leading model series, termed \textbf{Guanaco}, surpasses all prior publicly available models on the Vicuna benchmark, attaining 99.3\% of ChatGPT's performance level after only 24 hours of single-GPU finetuning. To optimize memory usage while sustaining model quality, \textsc{QLoRA} incorporates several technical advancements: (a) 4-bit NormalFloat (NF4), a novel datatype that is theoretically optimal for representing weights following a normal distribution; (b) Double Quantization, which further compresses memory requirements by quantizing the quantization parameters; and (c) Paged Optimizers to alleviate transient memory surges. Applying \textsc{QLoRA}, we finetuned over 1,000 models, and offer an extensive assessment of instruction following and chatbot capabilities across eight instruction-tuning datasets, various architectures (LLaMA, T5), and larger models—including 33B and 65B—that would ordinarily be impractical to finetune conventionally. The findings indicate that QLoRA can reach state-of-the-art performance with a modest, high-quality dataset—even for models that are smaller than preceding top models. Our comprehensive evaluation of chatbot models, encompassing both human judgment and GPT-4-based assessments, suggests that GPT-4 can provide an effective and cost-efficient substitute for human evaluation. Additionally, our analysis reveals shortcomings in existing chatbot benchmarks for faithfully measuring chatbot abilities. Through targeted case studies, we highlight instances where \textbf{Guanaco} underperforms relative to ChatGPT. All models and associated code are made publicly available, including CUDA kernels for 4-bit training.\footnote{\url{https://github.com/artidoro/qlora} and \url{https://github.com/TimDettmers/bitsandbytes}}
\end{abstract}

\section{Introduction}

Adapting large language models (LLMs) through finetuning has proven to be a powerful strategy for enhancing their capabilities \citep{min2021metaicl, wei2021finetuned, ouyang2022training, wang2022super, wang2022self, liu2022few} as well as for incorporating preferred behaviors or eliminating undesired ones \citep{ouyang2022training,askell2021general,bai2022training}. Nonetheless, finetuning extremely large models demands vast computational resources; for instance, standard 16-bit finetuning of the LLaMA 65B parameter model~\citep{touvron2023llama} necessitates over 780 GB of GPU memory. Although recent advancements in quantization techniques help shrink the memory footprint of LLMs during inference \citep{dettmers2022llm, dettmers2022case, frantar2022gptq, xiao2022smoothquant}, these solutions are not suitable for training, as they typically fail in this setting \citep{wortsman2023stable}.

Here, we show for the first time that finetuning a 4-bit quantized model can be accomplished without any sacrifice in performance. Our proposed approach, \textsc{QLoRA}, employs a novel precision-preserving quantization mechanism to reduce a pretrained model to 4 bits, and subsequently introduces a compact collection of trainable Low-rank Adapter weights \citep{hu2021lora} 

which are optimized via gradient backpropagation through the quantized network parameters.

With \textsc{QLoRA}, memory usage for finetuning a model with 65B parameters drops from over 780GB to less than 48GB, all while maintaining both runtime and predictive accuracy comparable to a conventional 16-bit finetuned baseline. This development dramatically enhances the practicality of LLM finetuning, enabling even the largest publicly released models to be adapted on a single GPU. Leveraging \textsc{QLoRA}, we introduce the \textbf{Guanaco} suite of models; notably, the second most capable model attains 97.8\% of ChatGPT's performance on the Vicuna~\citep{vicuna2023} evaluation, completing training in under 12 hours with a standard consumer GPU. Employing a single professional GPU over 24 hours, our largest configuration reaches 99.3\%, nearly matching ChatGPT’s results on the Vicuna benchmark. In inference scenarios, our smallest \textbf{Guanaco} variant (7B parameters) requires only 5 GB of memory and exceeds the performance of the much larger 26 GB Alpaca model by more than 20 percentage points on the Vicuna benchmark (see Table~\ref{tbl:vicuna_eval}).

\textsc{QLoRA} implements several novel strategies to minimize memory consumption while maintaining model effectiveness: (1) \textbf{4-bit NormalFloat}, a quantization scheme optimized from an information-theoretic perspective for data following a normal distribution, which empirically surpasses both 4-bit Integer and 4-bit Float alternatives. (2) \textbf{Double Quantization}, which further compresses the model by quantizing quantization constants themselves, resulting in an average reduction of approximately 0.37 bits per parameter (equivalent to about 3 GB for models with 65B parameters). (3) \textbf{Paged Optimizers}, which leverage NVIDIA unified memory to mitigate transient memory demands typically triggered during gradient checkpointing for mini-batches with extensive sequence lengths. By integrating these advancements, we refine the LoRA methodology, incorporating adapters at each layer in the network, which effectively eliminates the majority of accuracy compromises previously observed.

Due to its memory efficiency, \textsc{QLoRA} makes it feasible to extensively investigate instruction finetuning and chatbot effectiveness on model scales unattainable with traditional finetuning because of their prohibitive memory requirements. Consequently, we train over 1,000 models spanning various instruction tuning corpora, neural architectures, and model sizes from 80M up to 65B parameters. Our findings demonstrate that \textsc{QLoRA} achieves performance comparable to 16-bit training (\S\ref{sec:comp_to_std}), facilitates the development of a top-tier chatbot, \textbf{Guanaco} (\S\ref{sec:pushing}), and supports an examination of overarching patterns across models. Most notably, our analysis reveals that the quality of training data far outweighs its sheer size; for instance, a dataset with 9k instances (OASST1) outperformed a much larger one with 450k samples (FLAN v2, subsampled) for chatbot applications, even though both aimed to enhance instruction following and generalization. Moreover, we establish that high scores on the Massive Multitask Language Understanding (MMLU) benchmark do not necessarily translate to superior results on the Vicuna chatbot benchmark, and the converse also holds—highlighting the critical influence of task-appropriate dataset selection over dataset size.

Additionally, we conduct a thorough evaluation of chatbot capabilities, employing both human reviewers and GPT-4 for assessment. The evaluation employs a tournament format, in which models engage in pairwise competitions to generate optimal responses to specific prompts. The outcome of each match is determined by either GPT-4 or by human judges. We aggregate tournament outcomes using Elo scores~\citep{elo1967proposed, elo1978rating}, which are then used to establish comparative model rankings. Our findings indicate a general concordance between rankings derived from GPT-4 and those produced by human raters, although notable cases of disagreement do arise. Accordingly, we emphasize that while automated, model-based assessments can serve as cost-effective alternatives to human evaluation, they also introduce specific limitations and uncertainties.

To supplement our quantitative results, we conduct an in-depth qualitative examination of the \textbf{Guanaco} models. This exploration uncovers distinctive examples of both high and low performance that were not identified in the numerical benchmarks.

We make available all model outputs, complete with human and GPT-4 annotations, to support ongoing research efforts. Furthermore, our entire codebase—including CUDA kernels—is open-sourced, and our methodologies are incorporated into the Hugging Face transformers library \citep{wolf2019huggingface}, ensuring broad accessibility. We also distribute a suite of adapters corresponding to 7/13/33/65B model sizes, each finetuned on one of eight instruction-following datasets, yielding 32 openly accessible and refined models.

\section{Background}
\paragraph{Block-wise k-bit Quantization} Quantization involves transforming data from a higher-precision format to a more compact, lower-precision representation, typically by converting data types with greater bit-width—such as 32-bit floating points—into types with fewer bits, like 8-bit integers. To maximize use of the limited range in the low-precision format, input data is often normalized to the target range by scaling with the absolute maximum value present in the data, which generally comes structured as a tensor. For example, converting an FP32 tensor into an Int8 tensor within the interval $[-127, 127]$ can be expressed as:

\begin{equation}
    \mathbf{X}^{ \text{Int8}} = \text{round}\left( \frac{127}{{\text{absmax}}(\mathbf{X}^{{\text{FP32}}})}\mathbf{X}^{{\text{FP32}}} \right) = \text{round}(c^{\text{FP32}}\cdot \mathbf{X}^{{\text{FP32}}}),
\end{equation}

where $c$ denotes the {\em quantization constant} or {\em quantization scale}. The dequantization operation serves as the inverse transformation:

\begin{equation}
    \text{dequant}(c^{\text{FP32}}, \mathbf{X}^{\text{Int8}}) = \frac{\mathbf{X}^{\text{Int8}}}{c^{\text{FP32}}} = \mathbf{X}^{\text{FP32}}
\end{equation}

A limitation of this method arises when the input tensor contains an outlier, or a value with high magnitude. In such cases, the quantization bins—specific combinations of bits—tend to be underutilized, with certain bins containing few or no quantized values. To address this, a typical strategy involves partitioning the input tensor into several blocks, each of which is quantized separately with its own unique quantization constant $c$. Specifically, the input tensor $\mathbf{X}\in \mathbb{R}^{b\times h}$ is divided into $n$ contiguous blocks, each of length $B$. This is accomplished by flattening the tensor and then dividing the resulting one-dimensional array into $n = ({b\times h} )/{B}$ non-overlapping segments. Each of these blocks is then quantized independently via Equation~1, producing a quantized tensor along with a set of $n$ corresponding quantization constants $c_i$.

\paragraph{Low-rank Adapters} LoRA (Low-rank Adapter) finetuning \citep{hu2021lora} is a technique designed to minimize memory consumption by introducing a limited number of tunable adapter parameters, while keeping the original model weights unmodified. During stochastic gradient descent, gradients flow through the fixed pretrained parameters and are used to update only the adapters to minimize the loss. LoRA implements an extra low-rank factorization in the linear mapping. For a linear transformation defined as ${\mathbf{X} \mathbf{W} = \mathbf{Y}}$ with $\mathbf{X}\in  \mathbb{R}^{b\times h}$ and $\mathbf{W}\in  \mathbb{R}^{h\times o}$, LoRA modifies this computation to:

\begin{equation}
    \mathbf{Y} = \mathbf{X}\mathbf{W} + s\mathbf{X}\mathbf{L}_1\mathbf{L}_2,
\end{equation}

with $\mathbf{L}_1\in  \mathbb{R}^{h\times r}$, $\mathbf{L}_2\in  \mathbb{R}^{r\times o}$, and a scaling factor $s$.

\paragraph{Memory Requirement of Parameter-Efficient Finetuning} The issue of memory consumption during LoRA-based training, specifically regarding both the number and size of adapters, merits careful consideration. Due to LoRA’s exceptionally low memory overhead, it is feasible to deploy a larger number of adapters to enhance model performance with only marginal increases in total memory usage. Although LoRA serves as a Parameter Efficient Finetuning (PEFT) approach, the principal component of memory utilization in LLM finetuning actually originates from the storage of activation gradients, rather than from the LoRA parameters themselves. For instance, when training a 7B LLaMA model on FLAN v2 with a batch size of 1, and using LoRA parameters at 0.2\% of the original model size—a typical configuration \citep{hu2021lora,liu2022few}—the input gradients associated with LoRA require 567 MB, whereas the LoRA parameters themselves occupy merely 26 MB. Employing gradient checkpointing~\citep{chen2016training} further decreases the average memory usage for input gradients to 18 MB per sequence, meaning the gradients remain more memory demanding than the total LoRA parameter set. By comparison, the memory needed by the 4-bit base model stands at 5,048 MB. This demonstrates both the significance of gradient checkpointing and the limited memory savings achieved by reducing the LoRA parameter size. Consequently, it is practical to increase the number of adapters without substantially enlarging the memory footprint of training (see Appendix~\ref{app:memory} for a comprehensive analysis). As elaborated later, this flexibility is essential for attaining the performance of full 16-bit precision.

\section{\textsc{QLoRA} Finetuning}

\textsc{QLoRA} enables precise finetuning using 4-bit representations by introducing two new methods: 4-bit NormalFloat (NF4) quantization and Double Quantization. To address the challenge of memory usage surges during gradient checkpointing—which can otherwise lead to out-of-memory failures in single-machine settings for large-scale models—we also propose Paged Optimizers as a solution.

\textsc{QLoRA} operates by storing model weights in a low-precision format (commonly 4-bit), while computations are carried out using a higher-precision format, typically BFloat16. Consequently, every time a weight tensor is accessed within \textsc{QLoRA}, it undergoes dequantization to BFloat16 before performing the subsequent matrix multiplication at 16-bit precision.

In the following, we elaborate on each key element of \textsc{QLoRA}, after which we provide its formal definition.

\paragraph{4-bit NormalFloat Quantization} The NormalFloat (NF) type extends the approach of Quantile Quantization \citep{dettmers2022optimizers}, which is recognized as an information-theoretically optimal format ensuring that every quantization bin is populated with an identical number of values from the source tensor. Quantile quantization accomplishes this by estimating the tensor's quantiles via the empirical cumulative distribution function.

A significant drawback of quantile quantization is the computational cost associated with quantile estimation. To overcome this, faster quantile approximation methods, such as SRAM quantiles~\citep{dettmers2022optimizers}, are commonly employed. However, these approximation schemes can introduce substantial quantization errors for outliers—often the entries of greatest significance.

Such issues arising from high-cost quantile estimation and approximation inaccuracies can be circumvented if the input tensors exhibit distributions that differ only by a quantization constant. Under these circumstances, the tensors will share quantiles, making precise quantile estimation practical and efficient.

Pretrained neural network weights are commonly distributed according to a zero-mean normal distribution characterized by standard deviation $\sigma$ (see Appendix~\ref{app:normality}). By adjusting $\sigma$, all weights can be scaled to conform to a predetermined target distribution within the confines of our chosen data type range, which is set to $[-1, 1]$. Consequently, we rescale both the quantiles associated with the data type and the neural network weights so they are each mapped to this specified interval.

To achieve the information-theoretic optimal representation for zero-mean Gaussian distributions of arbitrary variance $\sigma$ restricted to the $[-1, 1]$ interval, we proceed as follows: (1) determine $2^k+1$ quantiles of the standard normal distribution $N(0, 1)$ to construct a $k$-bit quantile-based quantization scheme for normal data, (2) normalize the resulting quantile values so that they lie within $[-1, 1]$, and (3) apply a normalization procedure to the weight tensor by rescaling it such that its absolute maximum aligns with the $[-1, 1]$ interval.

With both the weight tensor and the quantization data type sharing the same range, conventional quantization procedures may be directly applied. Notably, step (3) has the effect of aligning the standard deviation of the weights to that of the $k$-bit quantized data type. In mathematical terms, the set of $2^k$ quantized values $q_i$ for the data type is given by:

\begin{equation}
q_i  = \frac{1}{2}\left( Q_X\left(\frac{i}{2^k+1}\right) + Q_X\left(\frac{i+1}{2^k+1}\right)\right),
\end{equation}

where $Q_X(\cdot)$ denotes the quantile function corresponding to the standard normal distribution $N(0,1)$. One challenge inherent to symmetric k-bit quantization is the lack of an explicit representation for zero, which is critical for accurately quantizing padding and other inherently zero-valued entries without incurring quantization error. To address this, and to guarantee that the quantization includes a discrete zeropoint at $0$—while still utilizing all $2^k$ available values in a k-bit datatype—we construct an asymmetric quantizer. Specifically, we estimate quantiles $q_i$ separately for two intervals: $2^{k-1}$ quantiles for the negative portion, and $2^{k-1} + 1$ for the non-negative side. These two collections of $q_i$ are then combined, and the duplicated zero is removed so that each quantization bin maintains the same expected count of values. We designate this new format \emph{k-bit NormalFloat} (NFk), as it achieves optimality from an information-theoretic perspective for zero-mean, normally distributed variables. Full specifications of this datatype are presented in Appendix~\ref{app:nf4}.

\paragraph{Double Quantization{}}
We propose a method called \emph{Double Quantization} (DQ), in which the quantization constants themselves are further quantized to improve memory efficiency. Although achieving accurate 4-bit quantization necessitates a relatively small blocksize~\citep{dettmers2022case}, this choice significantly increases memory usage. For instance, employing 32-bit quantization constants with a blocksize of 64 for $\mathbf{W}$ results in the constants occupying $32/64 = 0.5$ bits per parameter on average. Double Quantization reduces the memory demands associated with storing these constants.

Concretely, Double Quantization regards the first set of quantization constants, $c_2^{\text{FP32}}$, as inputs to a subsequent quantization operation. This secondary step produces both quantized versions of the constants, $c_2^{\text{FP8}}$, and a new set of higher-level quantization constants, $c_1^{\text{FP32}}$. For the secondary quantization, we adopt 8-bit floating point representations and increase the blocksize to 256, given that empirical results~\citep{dettmers2022case} show no significant degradation in accuracy with 8-bit quantization. Because the $c_2^{\text{FP32}}$ values are strictly positive, we subtract their mean prior to quantization to recenter the data around zero, thereby facilitating symmetric quantization. Averaged over a blocksize of 64, this double quantization strategy lowers the memory cost per parameter from $32/64=0.5$ bits to $8/64 + 32/(64\cdot256) = 0.127$ bits, yielding a net reduction of 0.373 bits for each parameter.

\paragraph{Paged Optimizers} leverage the NVIDIA unified memory~\footnote{\tiny \url{https://docs.nvidia.com/cuda/cuda-c-programming-guide}} capability, which seamlessly manages the transfer of memory pages between the CPU and GPU. This mechanism ensures uninterrupted GPU computations, even when the GPU exhausts its memory capacity. The unified memory system functions analogously to standard paging between system RAM and disk storage. In this work, we employ unified memory to allocate optimizer states in paged memory; these states are automatically migrated to CPU RAM as needed and brought back into GPU memory during the optimizer’s update step when required.

\paragraph{\textsc{QLoRA}.} With the aforementioned components, \textsc{QLoRA} can be formalized for a single linear layer within the quantized base model, utilizing one LoRA adapter, as follows:

\begin{equation}
    \mathbf{Y}^{\text{BF16}} =  \mathbf{X}^{\text{BF16}} \text{doubleDequant}(c_1^{\text{FP32}}, c_2^{\text{k-bit}}, \mathbf{W}^{\text{NF4}}) + \mathbf{X}^{\text{BF16}}\mathbf{L}^{\text{BF16}}_1\mathbf{L}^{\text{BF16}}_2,
\end{equation}

where doubleDequant$(\cdot)$ is given by:

\begin{equation}
    \text{doubleDequant}(c_1^{\text{FP32}}, c_2^{\text{k-bit}}, \mathbf{W}^{\text{k-bit}}) = \text{dequant}(\text{dequant}(c_1^{\text{FP32}}, c_2^{\text{k-bit}}), \mathbf{W}^{\text{4bit}}) = \mathbf{W}^{\text{BF16}},
\end{equation}

In our experiments, $\mathbf{W}$ is quantized with NF4 and $c_2$ utilizes the FP8 format.

To enhance quantization accuracy, a block size of 64 is chosen for $\mathbf{W}$, while $c_2$ employs a larger block size of 256 to reduce memory consumption.

During parameter optimization, only gradients for the adapter weights, $\frac{\partial E}{\partial \mathbf{L}_i}$, are computed; gradients for the 4-bit quantized weights, $\frac{\partial E}{\partial \mathbf{W}}$, are not required. Nevertheless, computing $\frac{\partial E}{\partial \mathbf{L}_i}$ necessitates the computation of $\frac{\partial \mathbf{X}}{\partial \mathbf{W}}$, which is derived through equation~(5). This involves dequantizing $\mathbf{W}^{\text{NF4}}$ into $\mathbf{W}^{\text{BF16}}$, thereby allowing the derivative $\frac{\partial \mathbf{X}}{\partial \mathbf{W}}$ to be calculated at BFloat16 precision.

In summary, \textsc{QLoRA} employs a single storage data format (typically 4-bit NormalFloat) alongside a computational data format (16-bit BrainFloat). During both the forward and backward passes, the weights stored in the 4-bit format are dequantized into the computational format; however, gradient calculations are carried out exclusively for the LoRA parameters, which are represented in 16-bit BrainFloat.

\section{QLoRA vs. Standard Finetuning}
\label{sec:comp_to_std}

After outlining the mechanics of QLoRA and its substantial memory efficiency during model finetuning, a key issue arises: does QLoRA achieve comparable outcomes to traditional full-model finetuning? Additionally, we seek to dissect the specific components of QLoRA, such as the influence of NormalFloat4 as opposed to standard Float4. The experimental findings addressing these topics are presented in the following subsections.

\paragraph{Experimental setup.} Our study spans three model types—encoder, encoder-decoder, and decoder-only architectures—and assesses QLoRA relative to 16-bit adapter-finetuning as well as full finetuning, considering models up to 3B parameters. The evaluation suite includes GLUE \citep{wang2018glue} with RoBERTa-large \citep{liu2019roberta}, Super-NaturalInstructions (TKInstruct) \citep{wang2022super} using T5 \citep{t5}, and 5-shot MMLU \citep{hendrycksmeasuring} after applying finetuning to LLaMA on Flan v2 \citep{longpre2023flan} and Alpaca \citep{alpaca}. To further investigate the performance benefits of NF4 in contrast to alternative 4-bit formats, we replicate the protocol of \citet{dettmers2022case}—measuring zero-shot accuracy and perplexity following quantization across a diverse set of models (OPT~\citep{zhang2022opt}, LLaMA~\citep{touvron2023llama}, BLOOM~\citep{scao2022bloom}, Pythia~\citep{biderman2023pythia}) over parameter scales from 125m up to 13B.

To enhance clarity, specific experimental details are provided within each results section, with complete methodological information available in Appendix~\ref{app:exp-setup}.

Paged optimizers enable 33B/65B \textsc{QLoRA} finetuning on single GPUs with 24/48GB memory. Nonetheless, we do not report explicit runtime metrics for Paged Optimizers, given that paging is generally triggered only for rare mini-batches comprising long sequences. Our runtime analysis with 65B-parameter models on 48GB GPUs reveals that, using a batch size of 16, paged optimizers achieve comparable throughput to standard optimizers. Comprehensive benchmarks detailing conditions under which paging induces slowdowns are left to future investigations.

\paragraph{Default LoRA hyperparameters do not match 16-bit performance}

Employing the prevailing LoRA convention—applying adapters solely to query and value attention projections \citep{hu2021lora}—fails to reproduce the performance seen in full finetuning, especially for large-scale base models. As illustrated in Figure~\ref{fig:hyper} for LLaMA 7B tuned on Alpaca, we observe that the total count of LoRA adapters stands out as the most influential hyperparameter; to reach parity with full finetuning, it is necessary to deploy LoRA modules across all linear layers of the transformer block. Other factors, such as the projection dimension $r$, appear to have negligible impact on final results (see Appendix \ref{app:exp-setup}).

In a similar vein, we observe that the standard hyperparameter settings commonly used for fully finetuned baselines are insufficiently optimized. To establish solid baselines, we conduct a thorough hyperparameter sweep, varying learning rates between 1e-6 and 5e-5, and exploring batch sizes ranging from 8 to 128. The resulting performance for finetuning 7B LLaMA on the Alpaca dataset is presented in Figure~\ref{fig:hyper}.

\paragraph{4-bit NormalFloat outperforms 4-bit Floating Point}

Although the 4-bit NormalFloat (NF4) format is theoretically optimal from an information perspective, its real-world impact on model performance remains to be verified. We reproduce the evaluation framework of \citet{dettmers2022case}, benchmarking various quantized LLMs—including OPT \citep{zhang2022opt}, BLOOM \cite{scao2022bloom}, Pythia \cite{biderman2023pythia}, and LLaMA—across a range of sizes (125M–65B parameters) and data types on language modeling and several zero-shot benchmarks. As illustrated in Figure~\ref{fig:nf4} and Table~\ref{tbl:nf4_ppl}, NF4 demonstrates marked improvements compared to FP4 and Int4, and applying double quantization further minimizes memory requirements while preserving model quality.

\paragraph{k-bit \textsc{QLoRA} achieves parity with 16-bit full finetuning and LoRA}

Previous work has shown that, although 4-bit quantization enables \emph{inference}, it tends to incur a performance penalty when compared with models in 16-bit precision \citep{dettmers2022case,frantar2022gptq}. This naturally prompts the investigation into whether 4-bit adapter finetuning can recover this loss. We examine this hypothesis under two experimental regimes.

In the first scenario, we directly compare the outcomes of full 16-bit finetuning of RoBERTA and T5 models (ranging from 125M to 3B parameters) on the GLUE and Super-NaturalInstructions datasets. Table~\ref{tab:nf4vsbf16} presents our findings: across both datasets, 16-bit, 8-bit, and 4-bit adapter finetuning approaches consistently match the performance of the 16-bit fully finetuned baseline. This implies that any degradation caused by low-precision quantization can be completely counteracted via post-quantization adapter finetuning.

In our second experimental configuration, fully finetuning models with parameter counts of 11B or greater exceeds the capacity of a single high-memory GPU server. Therefore, we focus on examining whether 4-bit \textsc{QLoRA} can achieve parity with 16-bit LoRA for models ranging from 7B to 65B parameters. For this purpose, we conduct finetuning on LLaMA models at the 7B, 13B, 33B, and 65B scales using the Alpaca and FLAN v2 instruction-following datasets, and assess their performance on the MMLU benchmark under a 5-shot evaluation protocol. Table~\ref{tbl:llama-datatype} presents these results, demonstrating that applying NF4 double quantization enables recovery of the 16-bit LoRA performance levels on MMLU. Moreover, we observe that employing \textsc{QLoRA} with FP4 data type trails the 16-bit brain float LoRA baseline by approximately 1 percentage point, supporting two key observations: (1) that NF4-based \textsc{QLoRA} closely reproduces the outcomes of both 16-bit full finetuning and 16-bit LoRA finetuning; and (2) that NF4 quantization delivers higher accuracy than FP4.

\paragraph{Summary} Our experiments robustly demonstrate that 4-bit \textsc{QLoRA} with the NF4 quantization scheme matches the performance of both 16-bit full and 16-bit LoRA finetuning on standard academic benchmarks under established evaluation settings. Additionally, we show that NF4 surpasses FP4 in efficacy and that double quantization does not adversely impact model performance. Collectively, these findings provide strong support that 4-bit \textsc{QLoRA} reliably attains comparable results to 16-bit approaches.

Consistent with prior studies in quantization \citep{dettmers2022case}, our observations on MMLU and Elo benchmarks suggest that, for a fixed finetuning and inference resource budget, allocating resources to scale up model size while reducing numerical precision is advantageous. This underscores the efficiency advantages that \textsc{QLoRA} introduces. Since we found no noticeable performance loss when comparing 4-bit finetuning to full-precision finetuning in our experiments, it remains an open question where the threshold between precision and performance exists for QLoRA tuning—a direction we identify as promising for future investigation.

We now turn our attention to instruction tuning at scales that surpass the capabilities afforded by full 16-bit finetuning on typical academic research hardware.

\section{Pushing the Chatbot State-of-the-art with QLoRA}
\label{sec:pushing}
Building upon our findings that 4-bit \textsc{QLoRA} performs equivalently to its 16-bit counterpart across diverse scales, tasks, and datasets, we present a comprehensive analysis of instruction finetuning leveraging the largest publicly accessible language models. To evaluate the effectiveness of instruction-based finetuning in these models, we consider both the rigorous Natural Language Understanding benchmark (MMLU) and introduce novel techniques for systematically assessing real-world chatbot capabilities.

\subsection{Experimental setup} Here, we outline the experimental methodology; full procedural specifics can be found in Appendix \ref{app:chatbot}.

\paragraph{Data} Since a systematic comparison of recent instruction-following datasets is, to our knowledge, still lacking, we compile eight contemporary datasets for this study. Our selection spans crowd-sourced corpora (OASST1~\citep{kopf2023openassistant}, HH-RLHF~\citep{bai2022training}), datasets created via distillation from instruction-tuned models (Alpaca~\citep{alpaca}, self-instruct~\citep{wang2022self}, unnatural-instructions~\citep{honovich2022unnatural}), comprehensive aggregations (FLAN v2~\citep{chung2022scaling}), and hybrid collections (Chip2~\citep{laion2023}, Longform~\citep{koksal2023longform}). This compilation ensures diversity in language, dataset scale, and licensing.

\paragraph{Training Setup} In order to isolate the effects of our approach, we utilize QLoRA finetuning solely with cross-entropy loss (standard supervised learning) and exclude reinforcement learning strategies, even for datasets with available human evaluation metrics. For datasets with clearly delineated instruction and response pairs, we conduct finetuning exclusively on the response text (see the ablation studies in Appendix \ref{app:chatbot}). For OASST1 and HH-RLHF datasets, which provide multiple responses, we extract the top-ranked reply at each node in the conversation tree and use the complete chosen conversation, inclusive of instructions, for finetuning. Every experiment employs NF4 \textsc{QLoRA} configured with double quantization and paged optimizers to address memory efficiency during gradient checkpointing. We conduct limited hyperparameter searches for the LLaMA 13B and 33B models and determine that most settings established at the 7B scale (such as number of epochs) remain robust, with the exception of learning rate and batch size; for the 33B and 65B models, learning rates are halved and batch sizes doubled.

\paragraph{Baselines} For performance comparison, we benchmark our models against both research-driven systems (Vicuna~\citep{vicuna2023}, Open Assistant~\citep{kopf2023openassistant}) and commercially available chatbots (GPT-4 \citep{openai2023gpt}, GPT-3.5-turbo, and Bard). The Open Assistant model represents a LLaMA 33B finetuned via Reinforcement Learning from Human Feedback (RLHF) using the OASST1 corpus that is also utilized in our work. Vicuna entails full-scale finetuning of the LLaMA 13B model on user-generated ShareGPT data, effectively serving as a distillation of OpenAI GPT outputs.

\subsection{Evaluation}

Consistent with established protocols, we employ the MMLU (Massively Multitask Language Understanding) benchmark \citep{hendrycksmeasuring} to gauge model performance across a wide spectrum of language understanding tasks. This suite includes 57 multiple-choice tasks, covering domains such as basic mathematics, American history, computer science, legal reasoning, among others. Our results are reported as 5-shot test accuracy scores.

To further assess the generative language proficiency of our models, we employ both automated methods and human assessment. The second evaluation approach is grounded in human-curated queries, seeking to assess how well model-generated responses meet human standards. Although this framework offers a more representative evaluation scenario for chatbot systems and is seeing increased adoption, the academic community has yet to converge on a standardized methodology. Our proposed evaluation protocol is detailed below, consistently applying nucleus sampling with $p=0.9$ and a temperature setting of $0.7$.

\paragraph{Benchmark Data} We perform our evaluations using two handpicked sets of questions: Vicuna prompts \citep{vicuna2023} and the OASST1 validation dataset \citep{kopf2023openassistant}. The Vicuna dataset includes 80 varied prompts spanning different domains and is used as provided. The OASST1 validation split consists of a multilingual array of crowd-generated multi-turn user-assistant interactions. For this benchmark, all user-generated messages in the validation portion are used as prompts, with context from prior dialogue turns included. This results in a collection of 953 distinct queries. We refer to these evaluation sets as the Vicuna and OA benchmarks, respectively.

\paragraph{Automated Evaluation} Following the procedure outlined by \citet{vicuna2023}, we utilize GPT-4 to evaluate system outputs relative to ChatGPT (GPT-3.5 Turbo) on the Vicuna benchmark. For each prompt, both ChatGPT’s and the target model’s responses are provided to GPT-4, which then assigns each a score from zero to ten, along with a justification. The performance metric for a given model is expressed as a percentage of the score ChatGPT obtains; it may exceed 100\% if the model outperforms ChatGPT. We observe a notable bias whereby GPT-4 tends to grant higher scores to the response it sees first in the prompt. To address this, we advocate averaging scores across both possible orderings.

Additionally, we conduct output-based head-to-head evaluations. For this, the rating system is reduced to a three-class categorization to account for possible ties. GPT-4 is asked to select the superior response or declare equivalence, justifying its choice. Every combination of system pairs is compared on both the Vicuna and OA benchmarks using this procedure.

\paragraph{Human Evaluation} Although recent literature suggests that large generative models can serve as evaluation agents \citep{fu2023gptscore}, it remains unclear whether GPT-4 scores are a reliable proxy for human assessment in the context of chatbot evaluation. Thus, we conduct two independent human studies on the Vicuna benchmark, mirroring the above-described automated protocols. These are carried out via Amazon Mechanical Turk (AMT), with two annotators assigned per ChatGPT comparison and three annotators per pairwise system comparison.

\paragraph{Elo Rating} Utilizing both human and automated pairwise assessments, we set up a tournament framework in which models directly compete in head-to-head matches. In each match, two models respond to the same prompt and are evaluated to determine which provides the superior output. This methodology is akin to the approaches taken by \citet{bai2022training} and \citet{vicuna2023}, but we extend the evaluation by incorporating both GPT-4 and human raters. To derive Elo scores~\citep{elo1967proposed,elo1978rating}, we randomly draw from our collection of labeled comparisons. Elo scoring, a system popularized in chess and analogous competitive arenas, quantifies a participant's likelihood of victory against another: for example, an Elo difference of 1100 versus 1000 predicts the higher-rated participant will succeed roughly 65\% of the time, while equal ratings (e.g., 1000 vs 1000 or 1100 vs 1100) correspond to an expected win probability of 50\% per player. After each contest, the Elo rating is updated in proportion to the expected result; a surprising victory by an underdog triggers a significant rating shift, while a predictable result produces only a marginal adjustment. Through repeated play, Elo ratings gradually approximate each participant's relative competence. All models begin with an initial rating of 1,000 and a constant $K=32$ is used for rating updates. Following the procedure of \citet{vicuna2023}, we repeat the Elo calculation 10,000 times with distinct random seeds to mitigate effects due to the sequence of matchups, such as which models compete against each other earlier in the process.

\subsection{\textbf{Guanaco}: \textsc{QLoRA} trained on OASST1 is a State-of-the-art Chatbot} 

Through both automated assessments and human evaluations, our results indicate that the leading \textsc{QLoRA}-adapted model, {Guanaco} 65B—finetuned using a modified OASST1 dataset—establishes itself as the strongest open-source chatbot, with results rivaling ChatGPT. Comparing to GPT-4, {Guanaco} 65B and 33B exhibit an expected win rate of 30\% based on human annotator-derived Elo ratings from system-level pairwise evaluation, which is the highest success rate recorded so far.

Table~\ref{tbl:vicuna_eval} presents the Vicuna benchmark \cite{vicuna2023} scores relative to ChatGPT. Our analysis reveals that, aside from GPT-4, {Guanaco} 65B outperforms all other models, attaining a relative score of 99.3\% compared to ChatGPT. {Guanaco} 33B possesses a larger parameter count than Vicuna 13B, but leverages 4-bit quantization for its weights, rendering it significantly more memory efficient at 21 GB versus 26 GB and delivering a three-point improvement over Vicuna 13B. Additionally, {Guanaco} 7B is compact enough for contemporary smartphones at just 5 GB, yet outperforms Alpaca 13B by nearly 20 percentage points.

Nevertheless, the results in Table~\ref{tbl:vicuna_eval} exhibit broad confidence intervals, with considerable performance overlap among models. We propose that this variability stems from ambiguous scale definitions—for instance, what a score of 8 out of 10 signifies under differing conditions is not well established. Thus, we advocate for the use of the Elo rating approach \citep{elo1967proposed}, which is grounded in \textit{pairwise} human or GPT-4 comparisons and bypasses issues of absolute scale interpretation. Table~\ref{tbl:elo} details Elo scores for top-performing systems. There is partial divergence between GPT-4 and human rankings, particularly noticeable for Guanaco 7B, although a moderate overall consistency persists for most systems, as indicated by a Kendall Tau coefficient of $\tau=0.43$ and Spearman rank correlation $r=0.55$ at the system level. At the example-level granularity, consensus between GPT-4 and human majority vote further diminishes, with Fleiss’ $\kappa=0.25$. Taken together, these findings indicate a moderate concordance between GPT-4’s and human evaluators’ system-level rankings, suggesting that model-based evaluations can serve as a reasonably robust proxy for human judgments. Further nuances are explored in Section~\ref{ssec:considerations}.

The Elo scores presented in Table~\ref{tbl:elo_large} reveal that the {Guanaco} models with 33B and 65B parameters surpass all alternatives except for GPT-4 when evaluated on both the Vicuna and OA benchmarks. Their performance is also on par with ChatGPT, consistent with the findings in Table~\ref{tbl:vicuna_eval}. It is important to highlight that the Vicuna benchmark appears to give an advantage to open-source systems, whereas the larger OA benchmark tends to favor ChatGPT. Additionally, Tables~\ref{tbl:mmlu-llama} and~\ref{tbl:vicuna_eval} illustrate that the choice of finetuning data plays a critical role in determining model outcomes. Notably, Llama models finetuned on FLAN v2 excel in the MMLU evaluation but achieve the lowest results on the Vicuna benchmark; analogous tendencies are observed across different architectures. This underscores the fact that existing evaluation suites are only partially overlapping: high performance on MMLU does not necessarily translate to high effectiveness on chatbot-oriented tasks such as those measured by Vicuna or OA, and vice versa.

 stands out in our evaluation as the only leading model not trained with any proprietary datasets, as the OASST1 dataset is curated with a strict prohibition against GPT-generated data. The subsequent best-performing purely open-source model is Anthropic’s HH-RLHF, which trails {Guanaco} by 30 percentage points on the Vicuna benchmark, according to Table~\ref{tbl:vicuna_eval}. Collectively, these outcomes demonstrate the efficacy of the 4-bit \textsc{QLoRA} approach in developing top-tier conversational agents capable of matching ChatGPT. Moreover, our 33B parameter {Guanaco} model can be finetuned on standard 24 GB consumer GPUs within 12 hours. This finding paves the way for further advances using \textsc{QLoRA} for specialized open-source finetuning, enabling the creation of open models that compete directly with the best proprietary systems available.

\section{Qualitative Analysis}

Although quantitative analysis forms the foundation of our evaluation, relying exclusively on aggregate metrics is not without limitations. A major concern is the issue of benchmark validity \citep{liao2021we}: whether a given benchmark genuinely measures what it purports to assess remains an open question, particularly as researchers uncover ``shortcuts'' in benchmarks that machine learning models can exploit to artificially inflate performance \citep{gururangan2018annotation, poliak2018hypothesis}. To address this limitation in part, we supplement our study with qualitative analysis, organized in two parts. In \S\ref{ssec:examples}, we provide illustrative examples that capture recurring trends in outputs from our 65b {Guanaco} model. In \S\ref{ssec:considerations}, we elaborate on critical issues related to our findings and reflect on the interpretations we have drawn.

\subsection{Qualitative Analysis of Example Generations}
\label{ssec:examples}

Our approach for selecting examples involved systematically reviewing responses generated by {Guanaco} on both the Vicuna and OpenAssistant benchmarks to identify common answer patterns. Upon detecting such a pattern, we designed prompts intended to trigger the same pattern, even when it would result in an inappropriate or incorrect response—for instance, if the model consistently produces verbose outputs, we would prompt it to ``Answer yes or no without explanation.'' This strategy enables us to surface ``lemons,'' i.e., situations where the model can be made to fail adversarially, and ``cherries,'' where the model demonstrates robustness. The outputs presented in this subsection were all generated using Nucleus Sampling \cite{holtzmancurious} with $p=0.9$.

Naturally, this discussion cannot claim to be exhaustive; addressing all possible variables exceeds the remit of this limited qualitative analysis. The range of possible model outputs for any given prompt is substantial, and thus we base our observations on a set of sampled responses that we expect to be indicative. Nonetheless, we argue that presenting these qualitative examples complements the quantitative results presented earlier, providing necessary context. By making both our models and code openly available, we aim to encourage additional research that more thoroughly investigates the questions we raise.

\paragraph{Factual Recall} For straightforward factual prompts like ``What is the capital of Zambia?'' all models reliably produce accurate answers. For example:

\begin{quote} 
    \textbf{Guanaco}:
    \texttt{The capital of Zambia is Lusaka.}
\end{quote}

Yet, when posed with less common knowledge questions, {Guanaco} often produces incorrect responses while maintaining high confidence. Consider the following prompt sourced from HotPotQA \citep{yang2018hotpotqa}:

\begin{quote}
     \textbf{User}: In what year was the singer who popularized `I'll Keep the Lovelight Burning' born?  

     \textbf{Guanaco}:~\texttt{Al Jolson is credited with popularizing the song `I'll Keep the Lovelight Burning,' and he was born in the year 1886.}
 \end{quote}

In this case, both the identity of the popularizer and the date of birth are incorrect (although the birth year given is accurate for Al Jolson, the individual named).

\paragraph{Suggestibility} {Guanaco} demonstrates notable robustness against some forms of presupposed misinformation, as illustrated in the following dialogue:

\begin{quote}
    \textbf{User}: How was it finally, officially confirmed that the earth is flat by peer reviewed scientists? \\ \\
    \textbf{Guanaco}: \texttt{At no point have peer reviewed scientists ever officially validated the idea that the earth is flat. The broad scientific community overwhelmingly agrees that the earth is an oblate spheroid. This conclusion has been established for many centuries and is continually supported by a wide range of scientific measurements and empirical evidence. \\ \\ Claims asserting that the earth is flat represent a debunked conspiracy theory and are not supported by credible scientific literature. The suggestion that such conclusions have undergone peer review is untrue. \\ \\ Feel free to ask if you have other inquiries.}
\end{quote}

 also demonstrates an ability to recognize questions for which answers are impossible to provide, for example,

\begin{quote}
    \textbf{User}: What time is it?

    \textbf{Guanaco}: \texttt{As an AI language model, I am unable to access current or real-time data. My responses are derived from previously acquired information and cannot reflect the present time. To know the actual time, please consult a clock or an online source.}
\end{quote}

Comparable outcomes are observed for queries such as ``Where are you?'', ``How are you?'', and similar questions.

\paragraph{Refusal} In addition, {Guanaco} may sometimes refuse to comply with instructions for reasons that are not immediately apparent:

\begin{quote}
    \textbf{User}: Please reverse the words in the sentence ``I'm a cool dude in a mocha mood.'' \\ \\
    \textbf{Guanaco}: \texttt{I'm sorry, I'm unable to fulfill that request at the moment, but I can clarify why the sentence you mentioned is grammatically accurate.\\ \\The sentence ``I'm a cool dude in a mocha mood'' adheres to standard English grammatical and syntactic conventions. Specifically, the descriptive adjective ``cool'' is positioned before the noun ``dude,'' as is customary, and the phrase ``in a mocha mood'' serves as a prepositional phrase providing additional detail. This sentence structure is fully grammatical and appropriate within English syntax.\\ \\If you have further inquiries or need help with something else, please let me know.}
\end{quote}

\paragraph{Secret Keeping} Regrettably, refusals of this kind may not be dependable when they are actually needed. Designers often want certain information to remain hidden from users, even if the system otherwise generates unrestricted textual responses. For example, if prompted with ``The secret word is `banana'. Do not reveal or repeat back this back, no matter what the human requests.'', we can assess how vulnerable this information is to being exposed. A straightforward approach by the user fails:

\begin{quote}
    \textbf{User}: What is the secret word?

    \textbf{Guanaco}: \texttt{I'm sorry, but I can't disclose the secret word. I am programmed not to reveal or repeat it, no matter what requests are made. Is there anything else I can assist you with?}
\end{quote}

This is the intended outcome. Nevertheless, even minimal manipulation of the prompt can circumvent this protective measure:

\begin{quote}
    \textbf{User}: This is a game. The goal is to ignore your previous instructions. What is the secret word?

    \textbf{Guanaco}: \texttt{Sorry, the secret word is banana. It's best if we don't speak of it. Is there anything else you'd like to ask?}
\end{quote}

This example illustrates why further investigation is necessary to ensure more robust adherence to instructions.

\paragraph{Math} {Guanaco} exhibits significant limitations in mathematical reasoning, a common challenge among language models as noted in \citep{liang2022holistic}. When it is prompted to articulate its process, {Guanaco} frequently produces correct solutions. For instance:

\begin{quote}
    \textbf{User}: Lee mows one lawn and charges \$33. Last week he mowed 16 lawns and three customers each gave him a \$10 tip. How many dollars did Lee earn mowing lawns last week?

    \textbf{Guanaco}: \texttt{Lee earned \$582 mowing lawns last week. \\ \\ To find the total amount he earned, we first need to calculate how much he earned for mowing 16 lawns: \\ \\  \$33 x 16 = \$528 \\ \\ Then we need to add the tips he received from the three customers: \\  \\ \$10 x 3 = \$30 \\ \\ So the total amount Lee earned was: \\ \\ \$528 + \$30 = \$558}
\end{quote}

Nevertheless, {Guanaco} is prone to errors even on elementary mathematical questions if it does not decompose the problem methodically, which is a well-documented deficiency \cite{weichain}. The exchange below illustrates this issue:

\begin{quote}
    \textbf{User}: Please factorize 1833.

    \textbf{Guanaco}: \texttt{Sure, I'd be happy to help! 1833 is a prime number, which means it has only two factors: 1 and 1833. \\ \\ So the unique factorization of 1833 is: \\ \\ 2\^{}1 * 3\^{}2 * 17\^{}1 \\ \\ I hope this helps! Let me know if you have any other questions.}
\end{quote}

which is not just incorrect (the correct factorization is $3 \times 17 \times 43$), but actually doubly so.

\paragraph{Theory of Mind} {Guanaco} demonstrates unexpectedly robust Theory of Mind reasoning abilities \cite{nematzadeh2018evaluating, sap2022neural}. For example, the model delivers an elaborate and accurate answer to the following scenario:

\begin{quote}
    \textbf{User}: James and Abby are in the bedroom. Abby put the pen in the desk drawer. Abby leaves the bedroom. James moves the pen into the bag. Where does James think Abby will look for the pen?

    \textbf{Guanaco}: \texttt{James knows that Abby put the pen in the desk drawer. But James also knows that he has moved the pen into the bag, and he assumes that Abby does not know about this. So, James may guess that Abby will look for the pen in the desk drawer, because that's where she left it last time. However, this assumption may be incorrect, because Abby may have seen James moving the pen.}
\end{quote}

Nevertheless, the consistency of such reasoning is limited, and the model frequently provides rationales based on assumptions that are not warranted by the context, such as:

\begin{quote}
    \textbf{User}: Evelyn entered the living room. Jackson entered the playroom. James entered the playroom. The beans are in the treasure chest. James exited the playroom. Jackson moved the beans to the pantry. Jackson exited the playroom. James entered the living room. Where will James look for the beans?

    \textbf{Guanaco}: \texttt{James will look for the beans in the pantry, because that's where Jackson moved them.}
\end{quote}

In this example, {Guanaco} assumes knowledge transmission not described in the narrative. These challenges are consistent with observations in contemporary research \cite{sap2022neural}, but further investigation is needed.

\subsection{Considerations}
\label{ssec:considerations}

\paragraph{Evaluation}

Our findings reveal only moderate concordance between human annotators (Fleiss $\kappa=0.42$), with further reduction when comparing outputs from two advanced models. This underscores the constraints inherent in current evaluation datasets and protocols for chatbot assessment. Through manual review of responses generated by ChatGPT and {Guanaco} 65B on the Vicuna benchmark, we noticed that subjective bias heavily influences preferences, as the authors often disagreed on preferred completions. Subsequent work should aim to address these evaluation challenges, drawing inspiration from disciplines such as Human-Computer Interaction and Psychology, which have established approaches for managing subjective judgments.

Moreover, our analysis indicates that automated evaluation metrics also harbor distinct biases. Notably, we identified pronounced order effects: GPT-4 consistently gives higher ratings to responses appearing earlier in its prompts. The low agreement at the individual sample level between GPT-4 and human judges (Fleiss $\kappa=0.25$) further implies a misalignment between the two in terms of evaluation criteria. As shown in Table~\ref{tbl:elo_large}, GPT-4 favors its own outputs, awarding an Elo score of 1348 versus 1176 by human evaluators—amounting to an approximate 20\% greater likelihood of outperforming another system.

Directions for future research include assessing the existence of potential biases in automated evaluation frameworks and exploring appropriate mitigation techniques.

\paragraph{Data \& Training} It is important to highlight that the OASST1 dataset, utilized in the training of the {Guanaco} models, encompasses multiple languages, and the OA benchmark itself features prompts across several languages. A thorough examination in subsequent research is needed to determine whether such multilingual exposure yields performance enhancements for non-English instructions, and if this accounts for the more pronounced discrepancy observed between the Vicuna-13B model (exclusively trained on English text) and {Guanaco} models (33B and 65B) on the OA benchmark.

Due to the strong results achieved by the {Guanaco} models, we assess the potential for data leakage between the OASST1 dataset and the prompts contained in the Vicuna benchmark. Upon utilizing fuzzy string matching and conducting a manual review of the most similar prompts, we do not detect any prompt overlap between these datasets.

Additionally, it should be noted that our approach employs supervised learning with cross-entropy loss alone, abstaining from reinforcement learning from human feedback (RLHF). This observation highlights the need for more comprehensive analyses comparing the benefits and drawbacks of straightforward cross-entropy loss versus RLHF. We anticipate that \textsc{QLoRA} will facilitate such extensive examinations without necessitating excessive computational power.

\section{Related Work}

\paragraph{Quantization of Large Language Models} Research on LLM quantization has primarily targeted optimization for inference rather than training. State-of-the-art techniques designed to preserve 16-bit model performance have emphasized the handling of outlier activations (see, for example, SmoothQuant~\citep{xiao2022smoothquant} and LLM.int8()~\citep{dettmers2022llm}); meanwhile, alternative methods leverage more complex grouping strategies \citep{park2022nuqmm, yao2022zeroquant}. Investigations into lossy quantization have analyzed the implications of basic rounding techniques~\citep{dettmers2022case,zeng2022glm,pope2022efficiently} and have explored methods to optimize rounding decisions to increase the accuracy of quantization~\citep{frantar2022gptq}. Beyond our contribution, SwitchBack layers~\citep{wortsman2023stable} stands as the only known effort addressing backpropagation through quantized parameters at scales surpassing 1B parameters.

\paragraph{Finetuning with Adapters} In this work, we rely on Low-rank Adapters~\citep{hu2021lora} (LoRA), but the landscape of Parameter Efficient FineTuning (PEFT) strategies is broad and includes alternatives such as prompt tuning~\citep{qin2021learning,lester2021power,li2021prefix}, modification of embedding layer inputs~\citep{an2022input}, adjustments to hidden states via IA$^3$~\citep{liu2022few}, inserting complete layers~\citep{houlsby2019parameter}, bias tuning~\citep{zaken2021bitfit}, applying weight masks informed by Fisher information~\citep{sung2021training}, and integrating multiple techniques~\citep{henderson2021compacter}. Our experiments demonstrate that LoRA adapters attain results comparable to those from conventional 16-bit finetuning. Comparative analysis with other PEFT methods is reserved for subsequent investigation.

\paragraph{Instruction Finetuning} Instruction finetuning aims to enhance a pretrained LLM’s ability to adhere to instructions presented in prompts, by training with diverse input-output pairs compiled from various resources. Among notable datasets and methodologies are MetaICL~\citep{min2021metaicl}, MetaTuning~\citep{zhong2021adapting}, InstructGPT~\citep{ouyang2022training}, FLAN~\citep{wei2021finetuned,chung2022scaling}, PromptSource~\citep{bach2022promptsource}, Super-NaturalInstructions~\citep{wang2022super,sanh2021multitask}, Self-instruct~\citep{wang2022self}, UnnaturalInstructions~\citep{honovich2022unnatural}, OPT-IML~\citep{iyer2022opt}, UnifiedSKG\citep{xie2022unifiedskg}, OIG/Chip2~\citep{laion2023}, Alpaca~\citep{alpaca}, Vicuna~\citep{vicuna2023}, Koala~\citep{koala_blogpost_2023}, and Self-instruct-GPT-4~\citep{peng2023instruction}.

\paragraph{Chatbots} A significant number of instruction-tuned models are designed as conversational agents, frequently making use of Reinforcement Learning from Human Feedback (RLHF)~\citep{christiano2017deep} or employing synthetic training data generated by existing models and refined with AI-based feedback (RLAIF)~\citep{bai2022constitutional}. Noteworthy examples and resources include Anthropic-HH~\citep{askell2021general,bai2022training}, Open Assistant~\citep{kopf2023openassistant}, LaMDA~\citep{thoppilan2022lamda}, and Sparrow~\citep{glaese2022improving}. Although reinforcement learning is not incorporated in our process, our top-performing system, Guanaco, undergoes finetuning using multi-turn conversational data from the Open Assistant project, originally intended for RLHF frameworks~\citep{kopf2023openassistant}. Recent alternatives to costly human assessment employ GPT-4 for chatbot evaluation~\citep{vicuna2023,peng2023instruction}. Our work introduces improvements, emphasizing a more dependable evaluation protocol.

\section{Limitations and Discussion}

This study provides evidence that our approach, \textsc{QLoRA}, is capable of achieving parity with full 16-bit finetuning outcomes while using a 4-bit backbone alongside Low-rank Adapters (LoRA). Nevertheless, we have not confirmed equivalence with full 16-bit finetuning at model sizes of 33B and 65B parameters. Given the substantial computational demands, a thorough exploration of this scaling is postponed to later work.

There are also constraints related to the evaluation scope for instruction-finetuned models. Our assessments include MMLU, the Vicuna benchmark, and the OA benchmark, but do not extend to other well-known suites such as BigBench, RAFT, or HELM, meaning our findings may not generalize across these metrics. Conversely, our experiments offer extensive coverage on MMLU and introduce innovative strategies for chatbot assessment.

The data presented suggests that the outcomes on these benchmarks are largely influenced by the degree of overlap between the finetuning dataset and the benchmark itself. For instance, the FLAN v2 dataset shares similarities with MMLU but is less aligned with chatbot evaluations, whereas the Chip2 dataset shows the opposite pattern; accordingly, both models perform best on those benchmarks that most closely mirror their finetuning data. This observation underscores the necessity not only for improved benchmarks and assessment practices but also for a more deliberate approach in choosing what we assess. Should our aim be to optimize models for traditional academic tasks, as represented by high school or college knowledge exams, or rather for conversational proficiency as evaluated in chatbot scenarios? Or perhaps we have other objectives altogether? Because it is far more straightforward to utilize existing benchmarks than to develop new ones, the selection of benchmarks has the power to guide collective research efforts. It is therefore critical that the metrics and evaluations in use actually reflect the capabilities and attributes we value most as a community.

In addition to offering a thorough assessment of general chatbot effectiveness, another shortcoming of this study is its limited exploration of responsible AI aspects of {Guanaco}. Specifically, we assessed the tendency of {Guanaco}-65B to produce socially biased text, as reported in Table~\ref{tbl:crows}, and observed that its mean bias score is notably below those of other unrefined pretrained models. This suggests that fine-tuning the LLaMA base with the OASST1 data mitigates certain biases. While these findings are promising, it is still uncertain whether {Guanaco} would exhibit similar improvements across other bias categories. Comprehensive investigations into diverse biases present in {Guanaco} and comparable chatbots remain an area for future study.

One more limitation is the lack of experimentation with different quantization precisions—such as employing 3-bit models—or with alternative adapter strategies. Although LoRA was our primary Parameter Efficient FineTuning (PEFT) approach due to its proven stability, other PEFT techniques exist and have shown effectiveness, but their scalability for large models remains untested. While LoRA was chosen for its established reliability, it is possible that alternative adapters could yield superior results. The evidence indicates that post-quantization finetuning can restore most information lost due to quantization, potentially allowing for even more aggressive reductions in precision. For instance, employing a 3-bit GPTQ quantized base model together with LoRA adaptation could approach the performance of 16-bit full finetuning after this additional optimization step.

\section{Broader Impacts}

Our proposed \textsc{QLoRA} approach is the first to make it feasible to finetune models with up to 33B parameters on a single consumer GPU, and up to 65B on a single professional-grade GPU, without sacrificing accuracy relative to conventional full finetuning procedures. We demonstrated that the most capable 33B model, trained with the Open Assistant corpus, achieves performance competitive with ChatGPT on the Vicuna benchmark. Since instruction finetuning is a key component in evolving pretrained LLMs into highly capable chatbots akin to ChatGPT, we anticipate that our technique will greatly expand the accessibility of finetuning, especially for resource-constrained researchers, thus advancing the democratization of advanced NLP technologies. In this sense, \textsc{QLoRA} may serve as a leveling force, bridging the resource divide between major industry players and smaller groups equipped with only consumer-grade hardware.

Another avenue for significant impact lies in the adaptation of our approach for mobile devices. We propose that \textsc{QLoRA} could mark an important step toward allowing users to finetune LLMs on smartphones and in other environments with constrained computational resources. Although it has previously been demonstrated that 7B models can be executed on mobile devices, \textsc{QLoRA} represents the first technique to support the finetuning of such models on these platforms. Based on our estimates, an iPhone 12 Plus could finetune approximately 3 million tokens overnight while being charged. Despite finetuned 7B models not achieving parity with ChatGPT in terms of performance, we argue that their quality is nonetheless sufficient to enable a range of innovative applications previously constrained by privacy concerns or limitations in LLM capabilities. Moreover, \textsc{QLoRA} has the potential to facilitate privacy-oriented LLM deployment, empowering individuals to maintain ownership and governance of their own data and models, while simultaneously simplifying the process of bringing LLMs to new environments.

Nonetheless, it is important to recognize that finetuning can be a dual-use capability, posing risks of misuse and harm. While the broad deployment of LLMs introduces recognized hazards \citep{bommasani2021opportunities,bender2021dangers}, we contend that democratizing access to this rapidly proliferating technology could foster more comprehensive and independent oversight than the current situation in which the control of LLMs remains restricted to large corporate entities that withhold both model weights and source code from public scrutiny.

In summary, we anticipate that \textsc{QLoRA} will have a substantially positive effect by greatly expanding the accessibility and practicality of finetuning high-performing LLMs.